%! Skills Focus: Selecting, Implementing, and Communicating Inference Procedures — Unit 7, Lesson 7.10 — Homework
\documentclass[10pt]{article}
\usepackage{apstats-article}
\usepackage{apstats-boxes}
\usepackage{grading}
\gradingassignment{apstats-u07-l10-hw}
\gradingcourse{advanced_statistics}
\newcommand{\TallMath}[1]{$\displaystyle #1\rule[-1.4em]{0pt}{3.2em}$}

\begin{document}

\pageheader{Unit 7, Lesson 7.10}{Homework}
\namedateperiod

\vspace{0.12in}

\begin{enumerate}

\item \textbf{Identify the Procedure.} \; For each scenario, name the correct inference
      procedure. No computation required.

\begin{enumerate}[label=\alph*.]
  \item A market researcher randomly selects 80 coffee drinkers and asks whether they prefer
        oat milk over dairy milk. She wants to estimate the proportion who prefer oat milk.

        Procedure: \makebox[9cm]{\smash{\gradedblank[id=hw1a, type=short, points=1, answer={{one-sample z-interval for p}}, check=text, display={\boldmath one-sample $z$-interval for $p$}]{8.8cm}{0.19in}}}

  \item A physical therapist records the number of days until pain relief for 12 randomly
        selected patients treated with a new therapy and 12 treated with a standard therapy.
        She wants to test whether the new therapy leads to faster relief, on average.

        Procedure: \makebox[9cm]{\smash{\gradedblank[id=hw1b, type=short, points=1, answer={{two-sample t-test for mu1 - mu2}}, check=text, display={\boldmath two-sample $t$-test for $\mu_1 - \mu_2$}]{8.8cm}{0.19in}}}

  \item A company tracks customer satisfaction scores (0--100) for 40 randomly selected
        customers before and after a new service policy. They want to test whether scores
        increased after the policy change.

        Procedure: \makebox[9cm]{\smash{\gradedblank[id=hw1c, type=short, points=1, answer={{matched-pairs t-test for mu_d}}, check=text, display={\boldmath matched-pairs $t$-test for $\mu_d$}]{8.8cm}{0.19in}}}
\end{enumerate}

\vspace{0.10in}

\item \textbf{Proportions: Hypothesis Test.} \;
      In a random sample of 120 city residents, 54 said they support expanding the public
      transit system. In a random sample of 100 suburban residents, 35 said the same.
      Is there evidence that the proportion who support expansion is higher among city
      residents? Use $\alpha = 0.05$.

      \textit{Write a complete four-step response.}

      \begin{enumerate}[label=\alph*.]
        \item \textbf{State:} Procedure: \makebox[7cm]{\smash{\gradedblank[id=hw2a1, type=short, points=1, answer={{two-sample z-test for p1 - p2}}, check=text, display={\boldmath two-sample $z$-test for $p_1 - p_2$}]{6.8cm}{0.19in}}}\\
              $H_0$: \makebox[5cm]{\smash{\gradedblank[id=hw2a2, type=short, points=1, answer={{p1 - p2 = 0}}, check=text, display={\boldmath$p_1 - p_2 = 0$}]{4.8cm}{0.19in}}} \quad $H_a$: \makebox[5cm]{\smash{\gradedblank[id=hw2a3, type=short, points=1, answer={{p1 - p2 > 0}}, check=text, display={\boldmath$p_1 - p_2 > 0$}]{4.8cm}{0.19in}}}
        \item \textbf{Plan:} Verify conditions (use numbers):
              \begin{itemize}
                \item Random: \gradedhide{\writeline}\gradedright{\gradedblank[id=hw2b1, type=rubric, points=1]{4.4in}{0.19in}}
                \item 10\%: \gradedhide{\writeline}\gradedright{\gradedblank[id=hw2b2, type=rubric, points=1]{4.4in}{0.19in}}
                \item Large Counts: \gradedhide{\writeline}\gradedright{\gradedblank[id=hw2b3, type=rubric, points=1]{4.4in}{0.19in}}
              \end{itemize}
        \item \textbf{Do:} $\hat p_1 =$ \makebox[1.5cm]{\smash{\gradedblank[id=hw2c1, type=short, points=1, answer=0.45, check=numeric, tolerance=0.005, display={\boldmath$0.45$}]{1.3cm}{0.19in}}} \; $\hat p_2 =$ \makebox[1.5cm]{\smash{\gradedblank[id=hw2c2, type=short, points=1, answer=0.35, check=numeric, tolerance=0.005, display={\boldmath$0.35$}]{1.3cm}{0.19in}}} \;
              $\hat p_c = \dfrac{54+35}{120+100} =$ \makebox[1.8cm]{\smash{\gradedblank[id=hw2c3, type=short, points=1, answer=0.405, check=numeric, tolerance=0.002, display={\boldmath$0.405$}]{1.6cm}{0.19in}}}

              $z = \dfrac{(\hat p_1 - \hat p_2)}{\sqrt{\hat p_c(1-\hat p_c)(1/120+1/100)}} \approx$ \makebox[1.8cm]{\smash{\gradedblank[id=hw2c4, type=short, points=1, answer=1.50, check=numeric, tolerance=0.03, display={\boldmath$1.50$}]{1.6cm}{0.19in}}} \quad $p$-value $\approx$ \makebox[2cm]{\smash{\gradedblank[id=hw2c5, type=short, points=1, answer=0.066, check=numeric, tolerance=0.004, display={\boldmath$0.066$}]{1.8cm}{0.19in}}}

        \item \textbf{Conclude:} \gradedhide{\writeline \writeline}\gradedright{\gradedblank[id=hw2d, type=rubric, points=2]{\linewidth-0.85in}{0.19in}}
      \end{enumerate}

\vspace{0.10in}

\item \textbf{Means: Confidence Interval.} \;
      A nutritionist randomly samples 22 adults and records their daily sodium intake.
      The sample mean is 2,580 mg and the sample standard deviation is 310 mg. A dotplot
      shows roughly symmetric data with one mild outlier. Construct a 95\% confidence interval
      for the mean daily sodium intake of all adults. Interpret your interval.

      \begin{enumerate}[label=\alph*.]
        \item \textbf{State:} Procedure: \makebox[7cm]{\smash{\gradedblank[id=hw3a1, type=short, points=1, answer={{one-sample t-interval for mu}}, check=text, display={\boldmath one-sample $t$-interval for $\mu$}]{6.8cm}{0.19in}}} \quad Parameter: \makebox[5cm][l]{\hspace{0.1cm}\smash{\gradedblank[id=hw3a2, type=rubric, points=1, display={\bfseries\boldmath $\mu$: mean sodium (mg)}]{4.3cm}{0.19in}}}
        \item \textbf{Plan:} Conditions:
              \begin{itemize}
                \item Random: \gradedhide{\writeline}\gradedright{\gradedblank[id=hw3b1, type=rubric, points=1]{4.4in}{0.19in}}
                \item Normal ($n < 30$): \gradedhide{\writeline}\gradedright{\gradedblank[id=hw3b2, type=rubric, points=1]{4.4in}{0.19in}}
              \end{itemize}
        \item \textbf{Do:} $t^* =$ \makebox[1.5cm]{\smash{\gradedblank[id=hw3c1, type=short, points=1, answer=2.080, check=numeric, tolerance=0.005, display={\boldmath$2.080$}]{1.3cm}{0.19in}}} (df = \makebox[1.5cm]{\smash{\gradedblank[id=hw3c2, type=short, points=1, answer=21, check=numeric, tolerance=0, display={\boldmath$21$}]{1.3cm}{0.19in}}}); \;
              $ME = t^* \cdot \dfrac{s}{\sqrt{n}} =$ \makebox[3cm]{\smash{\gradedblank[id=hw3c3, type=short, points=1, answer=137.4, check=numeric, tolerance=0.5, display={\boldmath$137.4$}]{2.8cm}{0.19in}}}; \; Interval: \makebox[4cm][l]{\hspace{0.1cm}\smash{\gradedblank[id=hw3c4, type=rubric, points=1, display={\bfseries\boldmath $(2442.6,\ 2717.4)$}]{3.3cm}{0.19in}}}
        \item \textbf{Conclude:} \gradedhide{\writeline \writeline}\gradedright{\gradedblank[id=hw3d, type=rubric, points=2]{\linewidth-0.85in}{0.19in}}
      \end{enumerate}

\vspace{0.10in}

\item \textbf{Matched Pairs.} \;
      A coach records the 40-meter sprint time (seconds) for 10 randomly selected sprinters
      before and after a 4-week speed training program. The differences (before $-$ after)
      have $\bar d = 0.31$ s and $s_d = 0.18$ s. The dotplot of differences is roughly
      symmetric with no outliers. Is there evidence that sprint times improved?

      Write a complete four-step response. Use $\alpha = 0.05$.

      \begin{enumerate}[label=\alph*.]
        \item \textbf{State:} Procedure: \makebox[7cm]{\smash{\gradedblank[id=hw4a1, type=short, points=1, answer={{matched-pairs t-test for mu_d}}, check=text, display={\boldmath matched-pairs $t$-test for $\mu_d$}]{6.8cm}{0.19in}}}\\
              $H_0$: \makebox[4.5cm]{\smash{\gradedblank[id=hw4a2, type=short, points=1, answer={{mu_d = 0}}, check=text, display={\boldmath$\mu_d = 0$}]{4.3cm}{0.19in}}} \quad $H_a$: \makebox[4.5cm]{\smash{\gradedblank[id=hw4a3, type=short, points=1, answer={{mu_d > 0}}, check=text, display={\boldmath$\mu_d > 0$}]{4.3cm}{0.19in}}}
        \item \textbf{Plan:} Conditions: Random: \gradedhide{\writeline}\gradedright{\gradedblank[id=hw4b1, type=rubric, points=1]{1.7in}{0.19in}}\quad Normal: \gradedhide{\writeline}\gradedright{\gradedblank[id=hw4b2, type=rubric, points=1]{1.5in}{0.19in}}
        \item \textbf{Do:} $t = \dfrac{\bar d - 0}{s_d / \sqrt{n}} = \dfrac{0.31}{0.18/\sqrt{10}} \approx$ \makebox[1.8cm]{\smash{\gradedblank[id=hw4c1, type=short, points=1, answer=5.45, check=numeric, tolerance=0.03, display={\boldmath$5.45$}]{1.6cm}{0.19in}}} \quad
              $p$-value $\approx$ \makebox[2cm]{\smash{\gradedblank[id=hw4c2, type=short, points=1, answer=0.0002, check=numeric, tolerance=0.0002, display={\boldmath$0.0002$}]{1.8cm}{0.19in}}} ($df =$ \makebox[1.5cm]{\smash{\gradedblank[id=hw4c3, type=short, points=1, answer=9, check=numeric, tolerance=0, display={\boldmath$9$}]{1.3cm}{0.19in}}})
        \item \textbf{Conclude:} \gradedhide{\writeline \writeline}\gradedright{\gradedblank[id=hw4d, type=rubric, points=2]{\linewidth-0.85in}{0.19in}}
      \end{enumerate}

\vspace{0.10in}

\item \textbf{AP-Style Multiple Choice.} \;
      \textit{A researcher wants to determine whether there is a difference in the mean exam
      scores of students who study alone versus those who study in groups. She randomly assigns
      30 students to study alone and 30 to study in groups, then records each student's exam
      score. Which procedure is most appropriate?}

      \begin{enumerate}[label=(\Alph*)]
        \item One-sample $t$-test for $\mu$
        \item Matched-pairs $t$-test for $\mu_d$
        \item Two-sample $t$-test for $\mu_1 - \mu_2$
        \item Two-sample $z$-test for $p_1 - p_2$
              \hfill\gradedright{\textit{Letter:}~\gradedblank[id=hw5, type=choice, points=1, answer=C]{0.4in}{0.2in}}
      \end{enumerate}

\end{enumerate}

\begin{extensionbox}
\textbf{Extension: CI vs.\ Significance Test} (optional)

A two-sample $t$-test for $\mu_1 - \mu_2$ gives $t = 2.31$, $p = 0.026$ at $\alpha = 0.05$.
A corresponding 95\% CI for $\mu_1 - \mu_2$ is $(0.4,\; 5.8)$.

Both lead to the same conclusion (reject $H_0$), yet the CI provides more information. Explain
in 2--3 sentences what the CI tells you that the significance test alone does not.
\end{extensionbox}

\begin{reflectionbox}
\textbf{Preview: Unit 8 --- Chi-Square Tests}

Unit 8 introduces a new class of inference procedures for categorical data with more than two
categories. Rather than a $z$- or $t$-statistic, these procedures use a chi-square statistic
($\chi^2$). The same four-step structure (State, Plan, Do, Conclude) will apply.

In the next unit, you will learn to test whether a categorical variable follows a claimed
distribution (goodness-of-fit), and whether two categorical variables are associated
(test of independence or homogeneity).
\end{reflectionbox}

\end{document}
